(1) 总质量为 \(M=\int_0^L\lambda(s)\,ds=\lambda_0L(1+\alpha/2)>0\)。左侧质量 \(M_{\rm l}(y)=\lambda_0[y+\alpha y^2/(2L)]\)，右侧为 \(M-M_{\rm l}\)，故净重为 \(F_g=g(2M_{\rm l}-M)\)。整个链所有质元速率相等，故动能严格为 \(K=M\dot y^2/2\)，并非 \(M_{\rm l}\dot y^2/2\)。虚功 \(\delta W=F_g\,\delta y\)，于是 \(M\ddot y=g[2\lambda_0y+\lambda_0\alpha y^2/L-M]\)。

(2) 从保守力积分 \(M\dot y\,d\dot y/dy=F_g(y)\) 得 \(\dot y^2=(2g/M)\{\lambda_0(y^2-y_0^2)+\lambda_0\alpha(y^3-y_0^3)/(3L)-M(y-y_0)\}\)。将右端记为 \(\Phi(y)\)，仅 \(\Phi(y)\ge0\) 的位置对给定初值可达。若初始静止但正处于平衡点，则需另外指定微小扰动才会产生非零速度。两端 \(y=0,L\) 是链完全脱离一侧的模型边界。

(3) 沿选定运动方向 \(dy/dt=\pm\sqrt{\Phi(y)}\)，所以从 \(y_0\) 到 \(y_1\) 的时间是 \(t=\left|\int_{y_0}^{y_1}dy/\sqrt{\Phi(y)}\right|\)，适用于其间无不可达点、按运动分支取号的情况。积分可能在平衡分界位置对数发散；对端点的平方根奇性则常为可积。左右段的质量持续交换，任何分别固定左右质量的牛顿质点模型都会漏掉穿过滑轮的动量通量；能量法自动保留了全链运动的惯性。

(4) 平衡条件为 \(M_{\rm l}(y_*)=M/2\)，即 \(y_*+\alpha y_*^2/(2L)=L(1+\alpha/2)/2\)。因 \(\lambda(s)>0\)，累积质量严格递增，故平衡唯一并位于内部。线性化 \(M\ddot{\eta}=2g\lambda(y_*)\eta\)，得 \(\eta\propto e^{\pm\gamma t}\)、\(\gamma=\sqrt{2g\lambda(y_*)/M}\)：这是势能极大值而非稳定振动。令 \(\alpha\to0\)，\(y_*=L/2\)、\(\gamma=\sqrt{2g/L}\)，与均匀柔链的经典不稳定模式一致。
